\documentclass[10pt,a4paper]{amsart}
\usepackage[margin=19mm]{geometry}
\usepackage{amsmath,amssymb,mathtools}
\usepackage[scr=rsfs,cal=euler]{mathalfa}
\usepackage{xfrac}
\usepackage[unicode,hidelinks,bookmarks=false]{hyperref}
\newcommand{\ie}{i.e.\ }
\newcommand{\cf}{cf.\ }
\newcommand{\resp}{respectively\ }
\newcommand{\Subsection}{\subsection}
\providecommand{\nobreakdash}{\nobreak\hbox{-}}
\setlength{\emergencystretch}{2em}
\multlinegap0pt
\usepackage{bm}
\usepackage{bbm}
\usepackage{dsfont}
\usepackage{xspace}
\usepackage{cancel}
\usepackage{mathtools}
\usepackage{xcolor}
\usepackage{amsthm}
\makeatletter
\@ifundefined{subsection}{\newtheorem{subsection}{Subsection}}{}
\@ifundefined{theorem}{\newtheorem{theorem}[subsection]{Theorem}}{}
\makeatother
\title[On the Probabilistic Approximation in Reproducing Kernel Hil]{On the Probabilistic Approximation in Reproducing Kernel Hilbert Spaces}
\author{Dongwei Chen, Kai-Hsiang Wang}
\date{Source: arXiv:2409.11679v2 (CC BY 4.0)}
\begin{document}
\maketitle

\noindent\emph{Source and licence.} On the Probabilistic Approximation in Reproducing Kernel Hilbert Spaces. Version arXiv:2409.11679v2 (CC BY 4.0), distributed under the Creative Commons Attribution 4.0 International licence, as recorded in the arXiv licence metadata for this exact version and shown on the abstract page https://arxiv.org/abs/2409.11679v2. Full rights notice: licenses/arxiv-2409.11679-CC-BY-4.0.txt.\par\smallskip

\noindent\textit{Editorial scope.} Counted unit: the Theorem whose source label is \texttt{pexistance} in the article (source0.tex lines 226--236, source label \texttt{pexistance}), delivered together with the article's own complete original proof of it (source0.tex lines 239--288, one \texttt{proof} environment of 50 lines). No mathematics is written, repaired or re-derived: statement and proof are the article's own text line for line. Required context retained uncounted: none. Every definition, display and label of those lines is retained unchanged. Omitted from this package and not redistributed: 1--225, 237--238, 289--1061. Nothing inside a retained excerpt is omitted or altered: every retained body is delivered exactly as the article's lines are written, so no omission receipt of any kind accompanies this package. Citations: the retained excerpts carry no citation command at all, so no bibliography entry of the article is transcribed and no reference is redistributed. Reference adaptation: none was needed and none was made; every reference inside the delivered unit resolves inside it, so no unresolved reference or citation is produced. Printed slips kept literal: no printed word, symbol, hypothesis or number of the counted statement or of its proof was altered. Layout and class: the article's own class is \texttt{amsart} with packages including fontenc, amsfonts, mathrsfs, mathtools, amssymb, hyperref, cleveref, amsmath, blkarray, booktabs, bigstrut, pdflscape, graphicx, dutchcal; the wrapper is the lane's pinned \texttt{amsart} editorial wrapper at 10pt on A4, which restates the theorem-like environments the retained text needs and adds the article's own macro definitions that the retained text needs (none), with extra wrapper packages (bm, bbm, dsfont, xspace, cancel, mathtools, xcolor). The delivered numbering is the unit's own. Assets: no retained span contains a figure environment, a graphics-inclusion command, a PDF-page inclusion, a table or a TikZ picture; no image file is copied into this package, the package declares no asset dependency and no image file is redistributed with it. Licence: version arXiv:2409.11679v2 of this article is distributed under the Creative Commons Attribution 4.0 International licence, as recorded in the arXiv licence metadata for this exact version and shown on the abstract page https://arxiv.org/abs/2409.11679v2. A full rights notice is delivered at licenses/arxiv-2409.11679-CC-BY-4.0.txt. Characters: every retained excerpt is pure ASCII, so the delivered engine, XeLaTeX with its default Latin Modern text font, reports no missing character.
\begin{theorem} \label{pexistance} 
  Let $\mathscr{H}(K)$ be an RKHS on $X$ with the feature map $\phi$.  
  Let $\mu \in \mathscr{P}(X)$ and $g \in L^p(X,\mu)$, where $1 \leq p<\infty$.
  Assume that $\phi: X \rightarrow \mathscr{H}(K)$ is a continuous $p$-frame for $\mathscr{H}(K)$ with respect to $\mu$. 
  Then the following problem 
    \begin{equation*}\label{muPNorm}
    \underset{f \in \mathscr{H}(K)}{\text{inf}} \ \int_{X} \vert f(x) - g(x) \vert ^p d\mu(x)
\end{equation*}
admits an optimizer $\hat{f} \in \mathscr{H}(K)$. Furthermore, if $p>1$, 
the optimizer is unique.
\end{theorem}

\begin{proof}
Since $g \in L^p(X,\mu)$ and $f=0 \in \mathscr{H}(K)$, we have 
 \begin{equation*}
   I:= \underset{f \in \mathscr{H}(K)}{\text{inf}} \ \int_{X} \vert f(x) - g(x) \vert ^p d\mu(x) \leq \int_{X} \vert g(x) \vert ^p d\mu(x)  < + \infty ,
\end{equation*}
i.e., the problem $I$ is bounded.
Let $\{f_i\}_{i=1}^{\infty}$ be a minimizing sequence. Then
\begin{equation*}
    \int_{X} \vert f_i(x) - g(x) \vert ^p d\mu(x) \rightarrow I.
\end{equation*}
Thus, there exists a number $0<N< \infty$ such that for each $i$ \textcolor{blue}{sufficiently large}, $ \Vert f_i - g \Vert_{L^p(X,\mu)} \leq  N$.
Then
\begin{equation*}
    \Vert f_i \Vert_{L^p(X,\mu)} \leq \Vert f_i - g \Vert_{L^p(X,\mu)} + \Vert g \Vert_{L^p(X,\mu)} \leq  N + \Vert g \Vert_{L^p(X,\mu)}, \ \text{for such $i$}.
\end{equation*}
On the other hand, since $\phi: X \rightarrow \mathscr{H}(K)$ is a continuous $p$--frame for $\mathscr{H}(K)$ with respect to $\mu$ with some lower frame bound $A>0$, then we have
\begin{equation*}
    A \Vert f_i \Vert_{\mathscr{H}(K)}^p \leq \int_{X} \vert f_i(x) \vert ^p d\mu(x) = \int_{X} \vert \langle f_i, \phi(x) \rangle_{\mathscr{H}(K)} \vert ^p d\mu(x), \;\text{for any $i$.}
\end{equation*}
Combining these results for the above sufficiently large $i$, we get
\begin{equation*}
     \Vert f_i \Vert_{\mathscr{H}(K)}^p \leq \frac{1}{A} \int_{X} \vert f_i(x) \vert ^p d\mu(x)  =  \frac{ \Vert f_i \Vert_{L^p(X,\mu)}^p}{A} \leq \frac{(N + \Vert g \Vert_{L^p(X,\mu)})^p}{A} < + \infty.
\end{equation*}
Thus  $\{f_i\}_{i=1}^{\infty}$ is a bounded sequence in $\mathscr{H}(K)$. Then $\{f_i\}_{i=1}^{\infty}$ has a weakly convergent subsequence $\{f_{i_k}\}_{k=1}^{\infty}$; i.e., there exists  $\hat{f} \in \mathscr{H}(K)$ such that for any $h \in \mathscr{H}(K)$, $  \langle f_{i_k}, h \rangle \rightarrow \langle \hat{f}, h \rangle \ \text{as} \ k \rightarrow \infty.$
\textcolor{blue}{For any $x \in X$}, taking $h = \phi(x)$, we get $f_{i_k}(x) \rightarrow \hat{f}(x)$.
Now by Fatou's lemma, we get 
\begin{equation*}
     \int_{X}  \underset{k \rightarrow \infty}{\text{lim inf}} \ \vert f_{i_k}(x) - g(x) \vert ^p d\mu(x) \leq  \underset{k \rightarrow \infty}{\text{lim inf}} \int_{X}  \vert f_{i_k}(x) - g(x) \vert ^p d\mu(x).
\end{equation*}
Using the pointwise convergence and that $\{f_i\}_{i=1}^{\infty}$ is minimizing, we obtain
\begin{equation*}
     \int_{X}  \vert \hat{f}(x) - g(x) \vert ^p d\mu(x) \leq  \underset{f \in \mathscr{H}(K)}{\text{inf}} \int_{X}  \vert f(x) - g(x) \vert ^p d\mu(x).
\end{equation*}
Therefore, $\hat{f}$ is an optimizer. 
\par
Next, we show that the optimizer is unique when $p>1$.
 Let $\hat{f_1}$ and $\hat{f_2}$ be optimizers. Then by Minkowski's inequality, we have
\begin{equation*}
    \| \frac{\hat{f_1}+\hat{f_2}}{2}-g \|_{L^p(X,\mu)} \leq \|\frac{\hat{f_1}-g}{2} \|_{L^p(X,\mu)} + \|\frac{\hat{f_2}-g}{2} \|_{L^p(X,\mu)}=I^\frac{1}{p}.
\end{equation*}
Since $I$ is the infimum and $\frac{\hat{f_1}+\hat{f_2}}{2} \in \mathscr{H}(K)$, the equality must hold and we infer the following two cases:
\begin{enumerate}
    \item $\hat{f_2}-g=0$ $\mu$-a.e.
    In this case, we have $I=0$ and $\hat{f_1}=g=\hat{f_2}$ $\mu$-a.e.

    \item $\hat{f_1}-g=\lambda(\hat{f_2}-g)$ $\mu$-a.e.\@ for some number $\lambda\geq 0$. Then $ \|\hat{f_1}-g \|_{L^p(X,\mu)} = \|\hat{f_2}-g \|_{L^p(X,\mu)} = I^{\frac{1}{p}}$ implies $\lambda=1$, and hence $\hat{f_1}=\hat{f_2}$ $\mu$-a.e. 
  %  It easily follows that $\lambda=1$ and hence $\hat{f_1}=\hat{f_2}$ $\mu$-a.e. 
\end{enumerate}
In either case, we conclude $\hat{f_1}=\hat{f_2}$ $\mu$-a.e. 
\end{proof}

\end{document}
